\documentclass[aps,prb,reprint,superscriptaddress,longbibliography,floatfix]{revtex4-2}

\usepackage{amsmath,amssymb,bm}
\usepackage{booktabs}
\usepackage{graphicx}
\usepackage{hyperref}
\usepackage{xcolor}

\input{generated/headline_values.tex}

\begin{document}

\title{Effective Central Charges across Clean, Disordered, and Monitored
Criticality: Benchmarks and an Exploratory Learning-Induced
Metal--Insulator Transition}

\author{Xu Tian}
\email{tianxu@westlake.edu.cn}
\affiliation{Department of Physics, School of Science, Westlake University,
Hangzhou 310030, China}
\author{Huidan Tan}
\affiliation{Department of Physics, School of Science, Westlake University,
Hangzhou 310030, China}

\date{\today}

\begin{abstract}
Universal finite-size terms provide stringent tests of numerical calculations at
two-dimensional critical points, but disorder and monitored dynamics complicate
both their interpretation and their uncertainty.  We establish a reproducible
validation hierarchy using three benchmarks and then apply it to an exploratory
learning-induced metal--insulator transition.  For the clean Ising model we
obtain $c=\CleanMCCharge$ from Monte Carlo data and
$c=\CleanExactCharge$ from the exact transfer matrix.  For the random-bond
Ising model on the Nishimori line and a weakly self-dual disordered model we
find effective central charges $c_{\rm eff}=\NishimoriCharge$ and
$\WeakCharge$, respectively.  At the candidate monitored transition
$\phi/\pi=\LearningCandidate$, entanglement and Casimir fits instead yield
$\LearningEntanglementCharge$ and $\LearningCasimirCharge$.  Their strong
disagreement, together with an unbracketed transition and unstable anisotropy
calibration, makes the monitored central-charge inference inconclusive.  The
comparison shows how validated benchmarks, resampling at the correct
independence level, and explicit claim gates prevent precise-looking fits from
being mistaken for universal results.
\end{abstract}

\maketitle

\section{Introduction}
\label{sec:introduction}

Continuous phase transitions are distinguished by the loss of a microscopic
length scale.  If a control parameter $t$ measures distance from criticality,
the correlation length diverges as
\begin{equation}
 \xi(t)\sim \xi_0 |t|^{-\nu}.
 \label{eq:correlation-length}
\end{equation}
At $t=0$, fluctuations occur on all scales between the lattice spacing and
the system size.  The renormalization group turns this observation into a
calculus: coarse graining generates a flow in the space of Hamiltonians,
critical theories are fixed points of that flow, and perturbations are
classified as relevant, irrelevant, or marginal according to whether they
grow, shrink, or remain finite under rescaling
\cite{WilsonKogut1974}.  Microscopic systems attracted to the same fixed point
share critical exponents and scaling functions and therefore form a
universality class.  This is why a lattice magnet, a fermion network, and a
monitored quantum circuit can meaningfully be compared even though their
microscopic variables are entirely different.

Two-dimensional criticality is unusually constrained.  Under standard
assumptions, scale invariance is enhanced to conformal invariance, and local
conformal transformations generate the infinite-dimensional Virasoro
algebra.  A conformal field theory is characterized not only by scaling
dimensions and operator-product coefficients but also by the central charge
$c$, which appears in the anomalous term of the stress-tensor algebra
\cite{BelavinPolyakovZamolodchikov1984,
DiFrancescoMathieuSenechal1997}.  The central charge is sometimes described as
a measure of low-energy degrees of freedom, but its most operational meaning
for numerical work is as the coefficient of the conformal anomaly.  Mapping
the plane to a cylinder changes the vacuum energy by an amount fixed by this
anomaly.  Consequently, a periodic cylinder of circumference $L$ has a
universal Casimir correction proportional to $c/L^2$
\cite{Cardy1984,BloeteCardyNightingale1986,Affleck1986}.  A subsystem of
length $\ell$ on a periodic quantum chain has a logarithmic entanglement
entropy with the same coefficient \cite{CalabreseCardy2004}.  These two
observables provide independent routes from raw numerical data to a
dimensionless universal number.

The coefficient is attractive precisely because it is restrictive.  A
critical temperature or coupling depends on microscopic conventions, whereas
$c$ should not.  Reproducing $c=1/2$ for the two-dimensional Ising fixed point,
whose thermodynamics is exactly known \cite{Onsager1944}, tests far more than
the visual linearity of a graph.  It simultaneously tests the normalization
of the partition function, the sign of the free energy, the factor relating a
strip Lyapunov exponent to an areal density, the cylinder boundary
conditions, and the finite-size extrapolation.  In Monte Carlo work it also
tests the integration of an energy observable from a known reference point
and the treatment of serial correlation.  Cluster algorithms such as the
Wolff update suppress critical slowing down \cite{Wolff1989}, but they do not
remove the obligation to identify independent streams and to propagate their
uncertainty through the nonlinear fit.

The apparent economy of the central charge can therefore be deceptive.
The universal term is subleading to the bulk free energy, and an error that is
small on the scale of the total free energy may be large on the scale of its
$L^{-2}$ coefficient.  Corrections from irrelevant operators can imitate a
change in slope over a short size range.  Measurements made along one Markov
chain are correlated, so their row count is not their effective sample size.
Quadrature errors in thermodynamic integration are systematic rather than
statistical.  Finally, choosing a fit window after seeing which one agrees
with the target converts a validation calculation into a selection-biased
one.  A credible result thus needs exact or microscopic oracles, nested
discretizations, independent stochastic replicas, residual diagnostics, and
predeclared size-window tests.  The purpose of our clean benchmark is to
exercise this entire chain before any disorder-specific interpretation is
introduced.

Quenched randomness changes both the fixed point and the statistical object.
The Harris criterion predicts when spatial fluctuations of a local critical
coupling destabilize a clean fixed point \cite{Harris1974}; subsequent field
theory shows how random fixed points and logarithmic corrections emerge in
two-dimensional Potts and Ising systems \cite{Ludwig1987}.  More basically,
one must distinguish quenched and annealed averages.  If a bond realization
$J$ is static on the equilibration time scale, the physical free energy is
$[\ln Z_J]_{\rm dis}$, whereas annealed disorder gives
$\ln[Z_J]_{\rm dis}$.  These operations are not interchangeable.  The replica
method represents the quenched logarithm by a derivative or $n\rightarrow0$
limit of $[Z^n]_{\rm dis}$, so different replica sectors can carry different
effective central charges.  In nonunitary or disordered theories the
finite-size coefficient is accordingly denoted $c_{\rm eff}$: it is a
well-defined universal amplitude for a specified averaged observable, not
automatically the central charge of a unitary Hilbert-space theory.

The random-bond Ising model on the Nishimori line is an especially stringent
test.  A relation between temperature and the bond distribution produces a
gauge identity that fixes otherwise nontrivial disorder averages
\cite{Nishimori1981}.  Duality and gauge structure locate a multicritical
region \cite{NishimoriNemoto2002}, while transfer-matrix, supersymmetric, and
network formulations expose its random fixed point
\cite{GruzbergReadLudwig2001,MerzChalker2002}.  Numerical estimates near
$c_{\rm eff}=0.464$ provide a reference for the ordinary quenched free energy
\cite{HoneckerJacobsenPiccoPujol2001}.  This qualifier is essential.  A
Born-weighted trajectory ensemble or a higher-replica moment need not measure
the same logarithm and may legitimately yield another coefficient.  In
addition, widths obtained from the same disorder realization are correlated.
A hierarchical or paired bootstrap must resample the realization as a whole;
resampling widths or transfer rows as if independent would manufacture
precision.  Bootstrap resampling is useful here not because it eliminates
modeling choices, but because it carries the empirically observed covariance
through the complete fitting procedure \cite{EfronTibshirani1993}.

Disordered fermions add a complementary viewpoint.  The
Chalker--Coddington construction expresses localization as propagation
through a network of scattering nodes \cite{ChalkerCoddington1988}, and
free-fermion mappings connect related networks to random-bond Ising transfer
matrices \cite{ChoFisher1997,MerzChalker2002}.  Anderson transitions separate
localized and extended single-particle states; their universality depends on
dimension and on the symmetry constraints of the Hamiltonian
\cite{EversMirlin2008}.  The Altland--Zirnbauer classification organizes
those constraints by time-reversal, particle-hole, and chiral symmetries
\cite{AltlandZirnbauer1997}.  Its tenfold-way formulation classifies
topological insulators and superconductors
\cite{SchnyderRyuFurusakiLudwig2008,RyuSchnyderFurusakiLudwig2010,
Kitaev2009}.  Symmetry class DIII, relevant below, has time-reversal symmetry
squaring to $-1$ together with particle-hole symmetry.  A DIII
metal--insulator transition can therefore have scaling behavior that cannot
be inferred by simply importing the clean Ising value.

Our third benchmark lies between the conventional quenched magnet and the
open monitored problem.  It propagates a fermionic Gaussian state while
sampling outcomes from the state-conditioned Born rule.  Gaussian covariance
matrices permit efficient simulation of fermionic linear optics
\cite{Bravyi2005}, but efficiency alone is not validation.  Sequential Born
sampling correlates outcomes through the evolving state; roundoff can also
drive a covariance matrix away from the fermionic purity and antisymmetry
manifold.  Weak self-duality supplies an independent physical audit: electric
and magnetic event densities must agree within their joint uncertainty even
though individual trajectories do not obey a pointwise equality.  We combine
this test with Gaussian invariants, effective sample size, residual trends,
and fit-window perturbations.  The resulting coefficient is again an
effective central charge of a precisely specified weighted ensemble.

Measurement-induced phenomena broaden the problem further.  Unitary dynamics
normally spreads entanglement, whereas local measurements remove quantum
information.  Competition between these processes can produce an
entanglement transition between volume-law and area-law trajectory phases
\cite{LiChenFisher2018,SkinnerRuhmanNahum2019}.  Purification dynamics gives a
related operational diagnostic \cite{GullansHuse2020}; mappings of random
circuits to statistical mechanics explain how measurement outcomes generate
replica interactions and nonunitary critical theories
\cite{BaoChoiAltman2020,JianYouVasseurLudwig2020}.  Numerical studies find
nontrivial critical scaling in Clifford and Haar-random circuits
\cite{LuntSzyniszewskiPal2019,
ZabaloGullansWilsonGopalakrishnanHusePixley2020}, while monitored free
fermions can display extended critical regimes and distinct routes to an area
law \cite{AlbertonBuchholdDiehl2021}.  Thus ``measurement-induced'' names a
mechanism, not one universal fixed point.

The learning-induced model considered here combines trajectory weighting,
free-fermion evolution, topology, and localization.  A control angle $\phi$
changes the learned or decohering channel, and a transfer diagnostic is used
to search for a class-DIII metal--insulator transition
\cite{WangVasseurTrebstLudwigZhu2025}.  Three logical steps must remain
separate.  First, the transition must be localized: data must identify both
sides of a crossing rather than merely a monotone trend.  Second, the
critical observable must be mapped to a conformal geometry.  A lattice
transfer direction and a spatial cut generally have different metric factors,
so an anisotropy parameter plays the role of a velocity.  Third, universal
estimators should agree after this calibration.  Entanglement and Casimir
amplitudes depend on different raw measurements; agreement is therefore a
powerful test of the assumed critical point and observable mapping.

These requirements motivate a distinction between parameter uncertainty and
claim uncertainty.  A regression can return a narrow conditional interval for
an amplitude at a chosen $\phi$, yet the scientific statement ``this is the
universal central charge at the transition'' can remain unsupported if
$\phi_c$ is not bracketed or the anisotropy drifts.  Increasing Monte Carlo
steps attacks variance at fixed model, size, and parameter.  It does not
generate the missing side of a phase boundary, enlarge the accessible
finite-size range, or correct a biased conversion factor.  We therefore use a
predeclared claim gate: a universal interpretation is released only if phase
localization, calibration stability, and independent-estimator agreement all
pass.  Failed gates retain the numerical amplitudes as diagnostics while
preventing a precise-looking but physically underdetermined number from being
published.

Here we organize four frozen calculations into a single ascending validation
hierarchy.  The clean square-lattice Ising model tests normalization,
Monte Carlo thermodynamic integration, and the finite-size coefficient
against an exact transfer matrix.  The Nishimori random-bond Ising model tests
ordinary quenched averaging, gauge identities, and realization-level
resampling.  A weakly self-dual Born ensemble tests Gaussian propagation,
state-conditioned sampling, self-duality, and common-disorder comparisons.
Only after those benchmarks do we analyze the exploratory learning-induced
DIII candidate.  The same philosophy is applied throughout: validate the
microscopic generator independently of the asymptotic fit, preserve the true
independence unit in uncertainty estimation, perturb the finite-size model,
and separate a stable computation from an allowed physical claim.
Figure~\ref{fig:01-workflow} summarizes the resulting hierarchy.

\begin{figure*}[t]
  \includegraphics[width=\textwidth]{generated/fig01-workflow.pdf}
  \caption{Validation hierarchy used throughout this work.  Exact and clean
  benchmarks validate normalization and scaling; disordered benchmarks add
  quenched averaging and common-disorder controls; the monitored calculation
  adds transition localization, anisotropy calibration, and agreement between
  independent universal estimators.}
  \label{fig:01-workflow}
\end{figure*}

The article has two complementary outcomes.  Quantitatively, all three
benchmarks reproduce their reference values within bootstrap uncertainty,
providing end-to-end evidence for the numerical mappings used in each
ensemble.  Methodologically, the monitored calculation demonstrates why
benchmark success must not be transferred automatically to a new observable.
Its putative DIII transition is not bracketed by the available scan, the
anisotropy calibration is unstable, and the entanglement and Casimir
conversions disagree.  We report both amplitudes because they diagnose what
the present calculation measures, but we do not average them or label either
one a universal central charge.

A useful way to formalize this distinction is to divide the error budget into
four layers.  Sampling error is the fluctuation remaining when the ensemble,
geometry, parameter, and estimator are fixed; it is the part conventionally
summarized by a standard error.  Discretization and finite-size error arise
because numerical integration uses a finite grid and simulations use finite
$L$; these errors are probed by nested grids, correction terms, and size cuts.
Calibration error enters when a measured lattice amplitude must be multiplied
by a velocity or anisotropy factor.  Model-selection error enters when the
candidate critical point, scaling window, or replica observable is not
uniquely identified.  These layers respond to different computational
remedies.  Longer runs reduce the first, denser grids and larger widths address
the second, dedicated geometries constrain the third, and broader phase scans
or additional observables resolve the fourth.  Reporting only the covariance
matrix of the final regression silently treats the other three as zero.

The hierarchy also clarifies the role of positive and negative controls.  An
exact positive control asks whether a known universal coefficient can be
recovered with the same conventions used later.  A microscopic control asks
whether the stochastic generator samples the intended distribution, for
example through the Nishimori energy identity or a Born-rule probability.
A negative control asks whether an analysis that should not show the target
effect remains null.  None of these controls proves the final universality
class by itself.  Instead, each removes a family of alternative explanations:
normalization mistakes, incorrect disorder frequencies, state-update drift,
or underestimated serial correlation.  This layered logic is more demanding
than obtaining an acceptable $\chi^2$, but it is also more informative when a
calculation fails.

Reproducibility is treated here as part of the physics rather than as an
afterthought.  Random streams are fixed independently of thread scheduling;
the generator, processed tables, and summaries are linked by immutable paths
and hashes; plotted points are read from those numeric artifacts rather than
from manually edited graphics; and headline values in the manuscript are
generated from validated adapters.  A reader should therefore be able to
trace a quoted interval back to its resampling unit and a withheld claim back
to the Boolean diagnostics that blocked it.  This provenance does not make a
finite-size ansatz correct, but it makes assumptions visible and prevents a
later presentation step from changing the statistical object.

The remainder of the paper is organized as follows.  Section
\ref{sec:models} defines the four ensembles and their finite-size observables.
Section \ref{sec:methods} explains the Rust simulation kernels, Xoshiro256++
stream construction, Python analysis, correlation-aware resampling, and
claim gates.  Section \ref{sec:benchmarks} presents the three validated
central-charge benchmarks.  Section \ref{sec:learning} analyzes the
learning-induced candidate and traces each failed inference gate to a
specific diagnostic.  Section \ref{sec:discussion} compares error mechanisms
and identifies the next decisive calculations.  The appendices derive the
observable mappings, specify the bootstrap, and provide a machine-readable
audit of the withheld claim.

\section{Models and universal observables}
\label{sec:models}

\subsection{Clean and random-bond Ising models}

For Ising spins $s_i=\pm1$ on a square lattice, the reduced Hamiltonian is
\begin{equation}
  -\beta H=\sum_{\langle ij\rangle}K_{ij}s_i s_j .
  \label{eq:ising-hamiltonian}
\end{equation}
The clean benchmark uses $K_{ij}=K_c=\tfrac12\ln(1+\sqrt2)$, whereas the
quenched benchmarks draw bond variables once per disorder realization and
hold them fixed during thermal sampling.  On the $\pm J$ Nishimori line, the
bond distribution and inverse temperature obey the Nishimori condition
\cite{Nishimori1981}.  The weakly self-dual ensemble uses the self-dual
parameter specified by the frozen production data.

For an infinitely long cylinder of circumference $L$, conformal invariance
gives
\begin{equation}
  f(L)=f_\infty-\frac{\pi c_{\mathrm{eff}}}{6L^2}
  +a_4L^{-4}+\cdots ,
  \label{eq:casimir}
\end{equation}
where $f$ is the dimensionless free-energy density.  In a unitary clean model
$c_{\rm eff}=c$; for a quenched random system it denotes the effective central
charge extracted from the disorder-averaged free energy.  The averaging order
is essential: the quenched observable is proportional to
$[\ln Z]_{\rm dis}$, not $\ln[Z]_{\rm dis}$.

\subsection{Monitored network model}

The learning-induced problem is represented by a disordered free-fermion
network in symmetry class DIII, with a measurement or decoherence angle
$\phi$ controlling the trajectory ensemble.  Its topological phases and
possible intervening metallic regime are diagnosed by transfer-matrix
localization observables and by the scale dependence of entanglement
\cite{AltlandZirnbauer1997,ChoFisher1997,
WangVasseurTrebstLudwigZhu2025}.  We analyze the frozen scan near
$\phi/\pi=\LearningCandidate$ rather than retuning the candidate after
inspecting its central-charge fits.

Two scaling relations are used.  A periodic one-dimensional cut of length
$L$ has
\begin{equation}
 S(\ell,L)=\frac{c_{\rm eff}}{3}
 \ln\!\left[\frac{L}{\pi}\sin\!\left(\frac{\pi\ell}{L}\right)\right]
 +s_0+\delta S ,
 \label{eq:entanglement}
\end{equation}
while the cylinder free energy follows Eq.~\eqref{eq:casimir} after converting
the lattice aspect ratio with an anisotropy factor $\alpha$.  Agreement of
these two estimates is a necessary internal consistency test; it is not
imposed as a fitting constraint.

\begin{table*}[t]
\caption{The four calculations and the additional controls required as model
complexity increases.}
\label{tab:01-models}
\begin{ruledtabular}
\begin{tabular}{lll}
Model & Universal probe & Principal control\\
\hline
Clean Ising & Casimir term & Exact transfer matrix\\
Nishimori RBIM & Quenched Casimir term & Disorder bootstrap\\
Weak self-dual & Quenched Casimir term & Common disorder\\
Monitored DIII & Entanglement/Casimir & Claim gates\\
\end{tabular}
\end{ruledtabular}
\end{table*}

\begin{table*}[t]
\caption{Notation used across the four calculations.  A symbol has the same
geometric meaning wherever possible, while the measured ensemble is stated
explicitly when it changes.}
\label{tab:notation}
\begin{ruledtabular}
\begin{tabular}{cl@{\qquad}cl}
Symbol & Meaning & Symbol & Meaning\\
\hline
$L$ & transverse circumference or chain length
& $M$ & transfer or imaginary-time length\\
$K$ & reduced Ising coupling
& $p$ & antiferromagnetic-bond probability\\
$\phi$ & monitoring/decoherence control angle
& $\alpha$ & spatial--temporal anisotropy factor\\
$f(L)$ & free-energy density on a cylinder
& $\gamma_1(L)$ & leading Lyapunov/free-energy rate\\
$S(\ell,L)$ & interval entanglement entropy
& $c$ & unitary conformal central charge\\
$c_{\rm eff}$ & ensemble-specific effective charge
& ${\rm CI}_{95\%}$ & central 95\% bootstrap interval
\end{tabular}
\end{ruledtabular}
\end{table*}

\section{Numerical methods and validation hierarchy}
\label{sec:methods}

\subsection{Monte Carlo implementation}

The Ising simulations are implemented in Rust.  Every pseudorandom stream uses
Xoshiro256++ \cite{BlackmanVigna2021}, with seeds assigned before parallel
execution so that scheduling cannot alter a realization.  Thermal updates use
the Wolff single-cluster algorithm \cite{Wolff1989}; measurements are separated
by completed cluster updates, and warmup samples are discarded before
accumulation.  Energy-like observables are formed with a Rao--Blackwell
estimator whenever the conditional bond expectation is available.  This
reduces thermal variance without changing the quenched sampling unit.

The simulation executable writes plain numerical records and metadata.  Python
is used only for validation, data reduction, nonlinear fits, bootstrap
resampling, and publication graphics.  This separation keeps the
high-throughput stochastic kernel in a compiled language while making the
statistical transformations directly inspectable.

\subsection{Statistical estimators}
\label{sec:statistical-estimators}

For each bootstrap replicate we resample independent units and repeat the full
fit, rather than perturbing the final parameter covariance matrix.  Clean
thermal streams are the independent units.  For random models a whole disorder
realization---including every width measured from it---is resampled as a
block.  This paired bootstrap \cite{EfronTibshirani1993} preserves
cross-width correlations.  When two nearby parameter values are compared, the
same bond realization and random labels are reused: this common-disorder
construction reduces the variance of their difference without pretending that
the pair is independent.

The reported standard error is the standard deviation of the bootstrap
parameter distribution, and the interval is its central 95\% percentile
interval.  Fit-window variation, truncation of small sizes, and correction
terms are treated as systematic checks rather than folded into a single
regression error.  For the exact clean calculation no stochastic error is
assigned.

\paragraph{Independence and uncertainty levels.}
The phrase \emph{independent sampling unit} refers to the object that can be
redrawn without reusing stochastic history.  It is a complete Rust stream for
the clean calculation, a complete bond realization for the Nishimori
calculation, and a complete state-conditioned trajectory stream for the Born
ensemble.  Thermal uncertainty describes fluctuations conditional on a fixed
Hamiltonian or disorder sample.  Quenched-disorder uncertainty describes the
variation between independently drawn static environments.  The latter
cannot be estimated by extending the transfer direction of one environment:
more rows refine its conditional Lyapunov exponent but do not add a new
disorder draw.

\paragraph{Conditional variance reduction.}
Rao--Blackwellization replaces a noisy event indicator by its exact
conditional expectation given the current state.  Common-disorder
comparisons similarly reuse the same random environment at two parameter
values so that shared fluctuations cancel in their difference.  Both
procedures reduce variance without changing the target expectation.  Their
uncertainty calculations must nevertheless preserve the pairing; treating
the two members as independent would discard the covariance that produces
the gain.

\subsection{Systematic errors and fit selection}
\label{sec:systematic-errors}

We separate four systematic categories from the bootstrap variance.
Finite-size truncation is tested by removing the smallest width and by adding
the leading symmetry-allowed correction term.  Thermodynamic-integration
discretization is tested on nested coupling grids.  Transition-location
uncertainty is controlled by demanding a two-sided phase bracket rather than
assigning an error bar to a monotone scan.  Anisotropy uncertainty is assessed
from alternative spatial and temporal windows, not only from the standard
error of one regression.  Fit windows and correction powers are specified
before comparison with the reference value; their shifts are reported as
diagnostics and are not silently absorbed into a statistical error bar.

These distinctions determine which additional calculation is useful.  More
Monte Carlo steps improve thermal precision and, up to autocorrelation,
effective sample size.  More independent streams improve uncertainty across
stochastic histories.  Larger system sizes constrain irrelevant corrections
and the asymptotic limit.  A wider parameter scan localizes a transition, and
new aspect ratios calibrate anisotropy.  Exchanging one remedy for another can
make an error bar smaller without making the physical inference more accurate.

\subsection{Predeclared claim gates}

We distinguish a numerical fit from a physical claim.  Benchmark values must
be stable to the prescribed size cuts and statistically compatible with their
reference targets.  A monitored universal-central-charge claim additionally
requires (i) a transition bracketed by independent phase diagnostics,
(ii) a stable space--time anisotropy calibration, and (iii) mutually
compatible entanglement and Casimir estimates.  Failing a gate does not erase
the computed result; it changes the allowed statement from a universal
estimate to a diagnostic conditional on the assumed candidate point.

\section{Benchmark central charges}
\label{sec:benchmarks}

\subsection{Clean Ising calibration}

The clean calculation uses periodic $L\times M$ tori with $M=8L$ and
$L=4,6,\ldots,20$.  We obtain the free energy at $K_c$ in two independent
ways.  The first evaluates the exact strip transfer matrix.  The second
integrates the measured energy from the known $K=0$ normalization,
\begin{equation}
 F(K_c)=-N\ln2+\int_0^{K_c}\langle H\rangle_K\,dK ,
 \label{eq:thermodynamic-integration}
\end{equation}
on a nested 129-point coupling grid.  The shift between 65- and 129-point
Simpson quadrature is $0.001955$ in the final central charge, below the
Monte Carlo standard error.  The cluster streams pass first/second-half and
between-replica tests with the predeclared $|z|<4$ threshold.

Figure~\ref{fig:02-clean} shows that the stochastic and exact free energies
share the same Casimir slope.  The primary $L_{\min}=6$ fit, including an
$L^{-4}$ correction, gives
\begin{align}
 c_{\rm MC}&=\CleanMCCharge\pm\CleanMCSE, \nonumber\\
 c_{\rm exact}&=\CleanExactCharge .
\end{align}
The Monte Carlo 95\% interval
$[\CleanCILow,\CleanCIHigh]$ contains $1/2$.  Diagnostic windows give
$c_{\rm MC}=0.499053$ for $L_{\min}=4$ and $0.490505$ for $L_{\min}=8$;
their increasing uncertainty and mutual compatibility show why the central
window is used without selecting it a posteriori.

\begin{figure*}[t]
  \includegraphics[width=\textwidth]{generated/fig02-clean.pdf}
  \caption{Clean Ising calibration.  (Left) Critical free-energy density
  versus $L^{-2}$ from the exact transfer matrix and thermodynamic-integration
  Monte Carlo calculation.  (Right) Independent estimates of $c$ compared
  with the exact Ising value.  Error bars are stream-level bootstrap
  uncertainties.}
  \label{fig:02-clean}
\end{figure*}

\subsection{Nishimori random-bond Ising model}

The $\pm J$ ensemble uses antiferromagnetic-bond probability
$p=0.1092212$ and
$K_N=\tfrac12\ln[(1-p)/p]=1.0493604763$.  For each of eight independent
replicas, a matrix-free transfer product is propagated through 4096 burn-in
rows and $2\,097\,152$ measured rows at widths
$L=4,6,8,10,12,14$.  Normalizing after every row prevents overflow and turns
the accumulated logarithmic norms into the quenched free-energy density.

The joint-width fit and hierarchical bootstrap in
Fig.~\ref{fig:03-nishimori} give
\begin{align}
 c_{\rm eff}&=\NishimoriCharge\pm\NishimoriSE,\nonumber\\
 {\rm CI}_{95\%}&=[\NishimoriCILow,\NishimoriCIHigh].
\end{align}
This interval contains the reference value $0.464$
\cite{HoneckerJacobsenPiccoPujol2001}.  Removing $L=4$ shifts the center from
$0.456484$ to $0.461475$; a paired bootstrap of the difference contains zero.
Independent microscopic audits also pass: the measured negative-bond fraction
has $z=-0.524$, and a centered coupling derivative differs from the exact
Nishimori energy identity by only $7.46\times10^{-5}$.

\begin{figure*}[t]
  \includegraphics[width=\textwidth]{generated/fig03-nishimori.pdf}
  \caption{Nishimori benchmark.  (Left) Quenched free-energy density and the
  $L^{-2}+L^{-4}$ fit.  (Center) Hierarchical bootstrap distribution of
  $c_{\rm eff}$.  (Right) Stability under the lower-width cut.  A disorder
  realization is resampled together at all widths.}
  \label{fig:03-nishimori}
\end{figure*}

It is important to state the replica convention.  Our value tests the
\emph{ordinary quenched} random-bond Ising free energy and therefore the
reference $c_{\rm eff}\simeq0.464$.  The often-quoted value near $0.522$
belongs to a different Born-weighted or higher-replica observable in monitored
constructions.  The two numbers are not competing estimates of the same
quantity; substituting the higher-replica target into this benchmark would
change the ensemble being measured.

\subsection{Weakly self-dual Born ensemble}

The third benchmark evolves a fermionic Gaussian covariance matrix at
$\theta=\pi/4$ and $\beta=\beta'=0.8813735870$.  Sequential, state-conditioned
Born sampling generates correlated outcomes according to
\begin{equation}
 P(s|\Gamma)=\frac{1+s\tanh(\beta)
 \langle i\gamma_a\gamma_b\rangle_\Gamma}{2}.
 \label{eq:born-rule}
\end{equation}
The leading Lyapunov or Shannon-free-energy rate is fitted over
$L=6,8,\ldots,32$ as
$\gamma_1(L)=f_\infty L-\pi c_{\rm eff}/(6L)+a/L^3$.
This Gaussian-state representation is consistent with standard efficient
fermionic simulation methods \cite{Bravyi2005}.

The result,
\begin{align}
 c_{\rm eff}&=\WeakCharge\pm\WeakSE,\nonumber\\
 {\rm CI}_{95\%}&=[\WeakCILow,\WeakCIHigh],
\end{align}
contains the reference $0.447$.  As shown in
Fig.~\ref{fig:04-weak}, the maximum studentized residual is $1.772$, its
correlation with $1/L$ is $-0.056$, the minimum effective sample size is
$4517$, and the electric--magnetic self-duality statistic is $z=-1.460$.
Required fit variants span $0.00553$, smaller than the declared $0.01$
systematic tolerance.

\begin{figure*}[t]
  \includegraphics[width=\textwidth]{generated/fig04-weak-self-dual.pdf}
  \caption{Weak self-dual benchmark.  The panels show the finite-size fit,
  studentized residuals, electric--magnetic self-duality audit, and
  autocorrelation/effective-sample-size convergence.  The residual and
  self-duality diagnostics are evaluated independently of the target value.}
  \label{fig:04-weak}
\end{figure*}

\begin{table*}[t]
\caption{Benchmark central charges.  Parentheses denote one bootstrap standard
error; intervals are central 95\% bootstrap intervals.  The exact clean result
has no stochastic uncertainty.}
\label{tab:02-benchmarks}
\begin{ruledtabular}
\begin{tabular}{lccccc}
Model & Widths & Estimate & 95\% interval & Reference & Outcome\\
\hline
Clean Ising, Monte Carlo & $4$--$20$ & $\CleanMCCharge(\CleanMCSE)$
 & $[\CleanCILow,\CleanCIHigh]$ & $0.500$ & Pass\\
Clean Ising, exact & $4$--$20$ & $\CleanExactCharge$
 & --- & $0.500$ & Pass\\
Nishimori RBIM & $4$--$14$ & $\NishimoriCharge(\NishimoriSE)$
 & $[\NishimoriCILow,\NishimoriCIHigh]$ & $0.464$ & Pass\\
Weak self-dual & $6$--$32$ & $\WeakCharge(\WeakSE)$
 & $[\WeakCILow,\WeakCIHigh]$ & $0.447$ & Pass\\
\end{tabular}
\end{ruledtabular}
\end{table*}

The compact comparison in Fig.~\ref{fig:05-benchmarks} makes the hierarchy
visible: all three stochastic intervals cover their predeclared references,
but the meaning of the coefficient changes from a unitary central charge to a
quenched or Born-weighted effective central charge.

\begin{figure*}[t]
  \centering
  \includegraphics[width=0.48\textwidth]{generated/fig05-benchmark-comparison.pdf}
  \caption{Cross-benchmark comparison.  Circles and horizontal bars show
  estimates and 95\% intervals; crosses are reference values.}
  \label{fig:05-benchmarks}
\end{figure*}

\section{Exploratory learning-induced metal--insulator transition}
\label{sec:learning}

\subsection{Phase localization and candidate selection}

Measurement-induced entanglement transitions can exhibit conformal scaling,
but localization, trajectory weighting, and nonlinear replica limits all
affect its interpretation \cite{SkinnerRuhmanNahum2019,LiChenFisher2018}.
We therefore first validate the scanning observable in an XY slice.  Its
crossing lies in $\phi/\pi\in[0.24,0.25]$, inside the reference window
$[0.20,0.28]$, as shown in Fig.~\ref{fig:06-scans}.  This confirms that the
scan reacts correctly in the calibration geometry.

The generic DIII scan is different.  Its diagnostic score increases
monotonically from $0.023$ at $\phi/\pi=0.16$ to $1.387$ at $0.32$, but the
available points do not exhibit both sides of a transition crossing.  The
transition is not bracketed.  The boundary value
$\phi/\pi=\LearningCandidate$ is retained as a preselected exploratory
candidate, with the unobserved upper side extending to $0.32$; it is not
promoted to a critical point by the subsequent fits.

\begin{figure*}[t]
  \includegraphics[width=\textwidth]{generated/fig06-phase-scans.pdf}
  \caption{Phase diagnostics.  (Left) The XY validation scan brackets its
  transition between $\phi/\pi=0.24$ and $0.25$.  (Right) The generic-DIII
  scan rises across the sampled interval but does not bracket a transition.
  The vertical marker at $0.30$ is therefore an exploratory candidate rather
  than a localized critical point.}
  \label{fig:06-scans}
\end{figure*}

\subsection{Entanglement estimate}

At each width $L=8,12,\ldots,32$, the interval entropy is fitted to the chord
coordinate in Eq.~\eqref{eq:entanglement}.  The width-dependent coefficients
are then extrapolated using a prespecified model set; model weights are carried
into the uncertainty rather than choosing the visually most favorable curve.
Figure~\ref{fig:07-entanglement} shows both stages.  The resulting conditional
estimate is
\begin{align}
 c_{\rm eff}^{(S)}
 &=\LearningEntanglementCharge\pm\LearningEntanglementSE,\nonumber\\
 {\rm CI}_{95\%}
 &=[\LearningEntanglementCILow,\LearningEntanglementCIHigh].
\end{align}
The fit has $\chi^2/{\rm dof}=0.370$ and remains available after removing the
smallest width, but the covariance condition number is
$6.9\times10^4$.  The wide interval correctly reflects noisy per-width
coefficients rather than the smoothness of a single $L=32$ chord fit.

\begin{figure*}[t]
  \includegraphics[width=\textwidth]{generated/fig07-entanglement.pdf}
  \caption{Entanglement analysis at $\phi/\pi=0.30$.  (Left) Entropy versus
  conformal chord coordinate for the available widths.  (Right) Extrapolation
  of the width-dependent effective coefficients.  Model-selection uncertainty
  is included in the reported interval.}
  \label{fig:07-entanglement}
\end{figure*}

\subsection{Casimir estimate and failed consistency gates}

The free-energy amplitude can be converted to a central charge only after
calibrating the spatial-to-temporal anisotropy.  The frozen calculation gives
$\alpha=\LearningAlpha$ with a nominal interval
$[0.02584,0.02956]$.  However, its relative spread across admissible windows
is $2.69$ and the spatial/temporal decay mismatch is $0.958$.  Thus the
anisotropy calibration is unstable even though a single-window regression
error is small.

Using the nominal $\alpha$ gives the conditional Casimir value
\begin{align}
 c_{\rm eff}^{(F)}
 &=\LearningCasimirCharge\pm\LearningCasimirSE,\nonumber\\
 {\rm CI}_{95\%}
 &=[\LearningCasimirCILow,\LearningCasimirCIHigh].
\end{align}
The difference
$|c_{\rm eff}^{(F)}-c_{\rm eff}^{(S)}|=9.519$ exceeds the combined 95\%
compatibility threshold $2.446$: the estimators disagree.  Figure
\ref{fig:08-learning} displays the apparently stable Casimir regression next
to the unstable calibration and the final inconsistency.  These diagnostics
show why increasing Monte Carlo steps alone would not resolve the present
problem.  More steps reduce sampling variance, whereas the leading failures
are an unbracketed parameter range, calibration drift, and estimator bias or
mismatch.

\begin{figure*}[t]
  \includegraphics[width=\textwidth]{generated/fig08-learning-diagnostics.pdf}
  \caption{Learning-induced-transition diagnostics.  (Left) Casimir-amplitude
  fit at the exploratory candidate.  (Center) anisotropy estimates under
  calibration variations.  (Right) entanglement and nominally converted
  Casimir estimates, whose 95\% intervals do not overlap.}
  \label{fig:08-learning}
\end{figure*}

\par\medskip
\noindent\begin{minipage}{\columnwidth}
\refstepcounter{table}\label{tab:03-learning}
\noindent\textbf{TABLE \Roman{table}.} Conditional monitored-model results at
$\phi/\pi=\LearningCandidate$.  ``No'' means that the quantity is retained as
a diagnostic but not published as a universal central charge.
\par\smallskip
\begin{ruledtabular}
\begin{tabular}{@{}lcc@{}}
Quantity & Value & Gate\\
\hline
XY bracket & $[0.24,0.25]$ & Pass\\
DIII bracket & None & Fail\\
$c_{\rm eff}^{(S)}$ & $\LearningEntanglementCharge(\LearningEntanglementSE)$
 & Available\\
$\alpha$ & $\LearningAlpha$ & Unstable\\
$c_{\rm eff}^{(F)}$ & $\LearningCasimirCharge(\LearningCasimirSE)$
 & Available\\
Estimator agreement & No & Fail\\
Published central charge & \LearningPublished & Withheld\\
\end{tabular}
\end{ruledtabular}
\end{minipage}
\par\medskip

Consequently, this pair of numbers in Table~\ref{tab:03-learning} does not
constitute a universal-central-charge estimate.  They are conditional diagnostics at an
exploratory parameter value.  This claim boundary is part of the result, not a
qualification added after observing disagreement.

\section{Cross-model discussion}
\label{sec:discussion}

\subsection{What the benchmark sequence validates}

The common analysis separates three questions: whether a numerical observable
is stable, whether its statistical error is correctly estimated, and whether
the observable supports the intended universal interpretation.  The clean
model tests normalization, thermodynamic integration, the sign and factor of
the Casimir coefficient, and finite-size corrections.  Exact agreement rules
out a broad class of coding errors that an internal goodness of fit cannot
detect.

The Nishimori calculation then changes the statistical unit.  Rows from one
transfer product are correlated, and widths sharing a disorder realization
are paired.  Resampling rows or widths independently would underestimate the
slope uncertainty.  The identity and bond-frequency audits test the disorder
generator separately from the fit.  The weak self-dual model adds
state-conditioned Born sampling: Gaussian invariant errors below
$2.6\times10^{-13}$ and the electric--magnetic check show that the measured
coefficient is not merely a stable output of a drifting covariance matrix.

\subsection{Error budget and parameter meaning}

Table~\ref{tab:04-errors} distinguishes error sources that are often collapsed
into one bar.  Statistical uncertainty is reduced by additional independent
streams or disorder realizations.  Autocorrelation is reduced by longer runs
only until the effective sample size becomes adequate.  Finite-size bias
requires larger $L$ or controlled correction terms.  Transition-location and
anisotropy errors require new parameter coverage and calibration geometries;
repeating the same point more accurately cannot remove them.

\begin{table*}[t]
\caption{Error mechanisms, controls, and effective remedies.}
\label{tab:04-errors}
\begin{ruledtabular}
\begin{tabular}{llll}
Source & Diagnostic & Present control & Remedy if failed\\
\hline
Thermal variance & replica/half-chain $z$ & Wolff streams, blocking
 & more independent clusters\\
Quenched variance & bootstrap spread & realization-level paired bootstrap
 & more disorder realizations\\
Autocorrelation & ESS & block analysis & longer streams or better updates\\
Finite-size bias & $L_{\min}$ drift/residuals & $L^{-4}$ or $L^{-3}$ terms
 & larger widths\\
Quadrature bias & nested-grid shift & 65/129-point comparison
 & denser coupling grid\\
Transition location & phase bracket & XY positive control
 & extend/refine DIII scan\\
Anisotropy & window spread & spatial/temporal comparison
 & dedicated aspect-ratio calibration\\
Estimator interpretation & $S$--$F$ difference & independent estimators
 & resolve mapping before universality claim\\
\end{tabular}
\end{ruledtabular}
\end{table*}

The large monitored values are therefore informative without being universal.
They may reflect evaluation away from criticality, an incorrect anisotropic
conversion, large finite-width corrections, or genuinely different replica
observables.  The present data do not discriminate among these mechanisms.
In particular, the success of a smooth Casimir fit cannot substitute for
agreement with entanglement or for a bracketed transition.

\subsection{Next decisive calculation}

The efficient next step is not uniform growth of the existing sample count.
First, extend the DIII scan beyond $0.32$ until the localization diagnostic
shows a two-sided crossing, then refine that interval using common random
numbers.  Second, calibrate $\alpha$ from several aspect ratios at the
bracketed point and require stability under width cuts.  Third, recompute both
universal estimators with larger widths and bootstrap their difference using
shared disorder.  Only if the difference interval contains zero should a
combined central-charge estimate be quoted.  This staged allocation directs
computation at the dominant uncertainty rather than at the smallest error bar.

\section{Conclusion}
\label{sec:conclusion}

We have placed four finite-size calculations in one PRB-style validation
framework.  Clean Ising thermodynamic integration yields
$c=\CleanMCCharge(\CleanMCSE)$ and agrees with the independent exact value
$\CleanExactCharge$.  Realization-level resampling gives
$c_{\rm eff}=\NishimoriCharge(\NishimoriSE)$ for the ordinary quenched
Nishimori model, and state-conditioned Born sampling gives
$c_{\rm eff}=\WeakCharge(\WeakSE)$ for the weakly self-dual benchmark.  Each
reference lies inside its 95\% interval and independent implementation gates
pass.

At the learning-induced metal--insulator-transition candidate, entanglement
and nominal Casimir conversions produce
$\LearningEntanglementCharge(\LearningEntanglementSE)$ and
$\LearningCasimirCharge(\LearningCasimirSE)$, respectively.  Because the DIII
transition is unbracketed, the anisotropy is unstable, and the estimators are
incompatible, \LearningPublished{} is correctly recorded as the publication
status of a central charge.  The defensible result is a reproducible failure
of the universality gates together with a concrete refinement path.  Across
all four models, the central lesson is that more Monte Carlo steps improve
accuracy only when sampling noise is the limiting error; they cannot repair an
unlocalized critical point or an uncontrolled observable mapping.

\appendix

\section{Fit models and window selection}
\label{app:fit-models}

All fit forms follow from an extensive cylinder free energy plus the leading
conformal Casimir term.  For clean Ising and the Nishimori free-energy density
we use
\begin{equation}
 f(L)=f_\infty-\frac{\pi c_{\rm eff}}{6L^2}
      +\frac{a_4}{L^4}.
\end{equation}
Periodic translation-invariant geometry excludes an $L^{-1}$ density term,
and the $L^{-4}$ contribution represents the leading retained irrelevant
correction.  The primary clean window begins at $L_{\min}=6$; $L_{\min}=4$
and 8 are diagnostic fits.  The primary Nishimori window includes $L=4$
because only six widths are available, while the paired $L_{\min}=6$ fit
measures sensitivity to that choice.

The weak-self-dual observable is stored as an extensive rate,
\begin{equation}
 \gamma_1(L)=f_\infty L-\frac{\pi c_{\rm eff}}{6L}
             +\frac{a_3}{L^3},
\end{equation}
which becomes the preceding density form after division by $L$.  Fits that
raise $L_{\min}$, double the block length, discard $L=30$, or extend burn-in
are paired with the primary data.  The largest center shift and the total
span of required variants are compared with predeclared tolerances.

For monitored entanglement, a linear chord fit is performed separately at
each $L$ before the resulting coefficients are extrapolated.  Keeping these
two stages separate prevents the many interval points at one width from
masquerading as independent evidence for the infinite-width limit.
Alternative extrapolation powers receive fixed model weights.  The Casimir
fit uses the measured $\gamma_1(L)$ form, but conversion of its amplitude to
$c_{\rm eff}$ is conditional on $\alpha$.  Hence a stable amplitude cannot
override a failed anisotropy gate.

Window variation is a sensitivity analysis, not permission to select the
answer nearest a reference.  A variant is accepted when residual structure,
precision, and predeclared stability criteria pass.  Bootstrap replicates
repeat the entire fit for the same window; paired replicates quantify the
difference between windows.  This procedure exposes both sampling dispersion
and finite-size-model dependence while keeping their meanings distinct.

\section{Observable mappings}
\label{app:mappings}

For a periodic cylinder with transfer direction $M\gg L$, write
$-\ln Z=M g(L)$ and $f(L)=g(L)/L$.  Conformal invariance gives
\begin{equation}
 g(L)=f_\infty L-\frac{\pi c}{6L}+O(L^{-3}),
\end{equation}
which is Eq.~\eqref{eq:casimir} after division by $L$.  This derivation fixes
both the sign and the factor of six.  The clean transfer matrix supplies an
independent microscopic mapping to the same coefficient.

For the $\pm J$ model, $g(L)$ is obtained from the long-product Lyapunov
exponent.  The physical quenched object is
\begin{equation}
 f_q(L)=-\lim_{M\rightarrow\infty}\frac{[\ln Z_{L,M}]_{\rm dis}}{LM}.
\end{equation}
The sign convention of the stored logarithmic norm is converted once, before
the finite-size fit.  For the Born ensemble, $\gamma_1(L)$ is extensive in
$L$, so the universal correction is $L^{-1}$ rather than $L^{-2}$ until the
rate is divided by the width.

For monitored entanglement, the coefficient of the chord logarithm is first
estimated separately at every width.  The reported
$c_{\rm eff}^{(S)}$ is the large-$L$ extrapolation of those coefficients.
The Casimir conversion contains the independently calibrated $\alpha$;
therefore its uncertainty and stability are logically prior to interpreting
$c_{\rm eff}^{(F)}$.

\section{Bootstrap construction}
\label{app:bootstrap}

Let $y_{rL}$ denote the estimate for realization or stream $r$ at width $L$.
A bootstrap replicate draws indices
$r_1^\ast,\ldots,r_R^\ast$ with replacement and forms the complete vector
\begin{equation}
 \bm{\bar y}^{\,\ast}
 =\frac{1}{R}\sum_{j=1}^{R}
 (y_{r_j^\ast L_1},\ldots,y_{r_j^\ast L_n}).
\end{equation}
The nonlinear finite-size fit is repeated on this vector.  All widths from one
quenched realization consequently enter together, preserving the empirical
cross-width covariance.  The clean analysis applies the same construction to
independent Monte Carlo streams.  The Nishimori and weak-self-dual production
analyses use 4000 replicates with fixed analysis seeds, but the simulation RNG
streams remain independent of those seeds.

For common-disorder comparisons, parameter values $a$ and $b$ share the same
bootstrap index in every replicate.  The sampled quantity is
$\Delta c^\ast=c_a^\ast-c_b^\ast$, not the difference of two independently
resampled fits.  This is the correct uncertainty for a paired fit-window or
parameter comparison.  Percentile intervals are read directly from the
replicate distribution; no Gaussian shape is assumed.

\section{Claim-gate audit}
\label{app:gates}

The frozen run contains \LearningTasks{} independent random streams and widths
$8,12,\ldots,32$, with elapsed time \LearningRuntime{} s.  Its summary artifact
has SHA-256 digest
\begin{widetext}
\begin{center}
\texttt{\LearningSummaryHash}.
\end{center}
\end{widetext}
The claim gate evaluates three Boolean conditions.  First, the DIII phase
diagnostic must supply a lower and an upper boundary; it supplies neither a
complete bracket nor a crossing in the available interval.  Second, the
anisotropy must be stable across prescribed windows; its relative spread is
$2.69$.  Third, the independent 95\% intervals for the entanglement and
Casimir conversions must be compatible; their center difference $9.519$
exceeds the combined threshold $2.446$.

All three failures are recorded in the source summary as
\texttt{diii\_transition\_not\_bracketed},
\texttt{anisotropy\_unstable}, and
\texttt{estimator\_disagreement}.  No post-processing rule overrides them.
The machine-readable publication flag is \LearningPublished.  This audit
ensures that a reader can reproduce not only the numerical values but also the
logical reason they are not presented as a central charge.

\subsection{Reproduction workflow}

The analysis uses immutable result directories containing configuration,
manifest, raw records, processed summaries, and plots.  Rust executables
generate stochastic records with Xoshiro256++ streams; Python scripts load
hash-gated artifacts, perform resampling and fitting, and build the figures.
The manuscript values are generated from those adapters rather than copied by
hand.  The complete workflow is invoked by the package Makefile, which runs
tests, regenerates figures and TeX macros, compiles the RevTeX source, and
verifies the resulting PDF.

\bibliography{references}

\end{document}
